% This is before

\title{GRACE: A Narrative-Centered Socio-Technical Ecosystem for Elder Wellbeing}
\author{Author,~\IEEEmembership{Member,~IEEE,}
        B.Author,~\IEEEmembership{Fellow,~OSA,}
        and~C.Author,~\IEEEmembership{Life~Fellow,~IEEE}% <-this % stops a space        
\thanks{Manuscript created October, 2020; This work was developed by the IEEE Publication Technology Department. This work is distributed under the \LaTeX \ Project Public License (LPPL) ( http://www.latex-project.org/ ) version 1.3. A copy of the LPPL, version 1.3, is included in the base \LaTeX \ documentation of all distributions of \LaTeX \ released 2003/12/01 or later. The opinions expressed here are entirely that of the author. No warranty is expressed or implied. User assumes all risk.}}

\markboth{Journal of \LaTeX\ Class Files,~Vol.~18, No.~9, September~2020}%
{How to Use the IEEEtran \LaTeX \ Templates}

\maketitle
\begin{abstract}
Population aging is increasing demand for digital systems that support older adults’ wellbeing. Most existing systems emphasize safety, monitoring, and service coordination, but underrepresent identity, meaning, contribution, and relationship quality. This paper proposes GRACE (Growing, Resourceful, Active, and Enriched Life), a narrative-centered socio-technical ecosystem for elder wellbeing. GRACE is designed as an intelligent multiagent community computing framework that supports older adults in maintaining relationships with the spiritual self, family and friends, working colleagues, communities, and the wider world, including markets and public services. The paper makes four contributions. First, it introduces a narrative-centered theory of elder wellbeing grounded in relationship quality and contribution exchange. Second, it proposes the GRACE Index, a daily measure of whether a person remains “in the zone” of narrative selfhood during ordinary and difficult periods. Third, it extends the framework with three innovations: daily hero’s-journey story creation, an AI-based fast state predictor, and a smart stimulus agent for anticipatory support. Fourth, it presents an agent-based simulation design to evaluate these innovations. We argue that elder-oriented AI should move beyond deficit-oriented monitoring toward socio-technical ecosystems that preserve dignity, agency, reciprocity, and life authorship.
\end{abstract}

\begin{IEEEkeywords}
aging, wellbeing, narrative identity, multiagent systems, human-centered AI, socio-technical systems, community computing, eldercare technology.
\end{IEEEkeywords}